\chapter*{Abstract}
\noindent The advent of smart-contract-enabled blockchains has paved the way for a new form of decentralized organization that can operate under autonomous and predictable rules encoded in transparent smart contracts. Since these Decentralized Autonomous Organizations (DAOs) are inexpensive to create and can effectively exist forever on the blockchain, it is vital to determine what distinguishes an “active” DAO from an “inactive” one. To address this, we developed a simplified theoretical framework grounded in financial, social, and development indicators. We then examined the 12 largest public blockchains to construct a novel dataset of 1,332 deployed DAOs, each manually classified according to these criteria, and made the dataset publicly available for future research. Using this dataset, we trained an Long Short-term Memory (LSTM) model that classifies DAO activity with 75\% accuracy, outperforming any individual heuristic.

\textit{\textbf{Keywords: } DAO, LSTM, Machine Learning, Organizational Activity}

